\section{Exact shared-query separations: policies, lower bounds, and verifiable certificates}
\label{sec:joint-separations}
\begin{theorem}[Feedback reduces mean joint width with two shared queries]
\label{thm:shared-prefix-separation}
Let $F=\{x\in\mathbb R_+^6:\sum_hx_h=1\}$, and allow at most two
queries from $P_kx=\sum_{h\leq k}x_h$, $k=1,\ldots,5$.  For the two
contrasts with coefficient vectors
\[
 c^{(1)}=(1/2,0,1,1/2,1/2,0),\qquad
 c^{(2)}=(2/3,2/3,1,1/3,0,1),
\]
both initial widths are one.  Let $J_Q(x)$ be the larger terminal sharp
width and let $R^J_{\mathcal C}=\inf_{P\in\mathcal C}\sup_{x\in F}
\mathbb E J_{Q_P(x)}(x)$.  Then
\[
 R^J_{\mathrm{randomized\ adaptive}}=\tfrac12
 <\tfrac23=R^J_{\mathrm{randomized\ batch}}.
\]
\end{theorem}
\begin{proof}
A batch partitions the six coordinates into segments.  Its scalar width
is the sum, over segments, of segment mass times coefficient range.
After querying boundary~3 with answer $t$, the worst remaining joint
widths for second queries~5 and~2 are, respectively,
\[
 G_{35}(t)=\max(t,1/3),\qquad
 G_{32}(t)=\max(1/2,1-t).
\]
Query~5 when $t\leq1/2$ and query~2 otherwise.  This guarantees $J\leq1/2$.
At the single feasible record $x^*=(0,0,1/2,0,0,1/2)$, all ten final
two-query sets have $J_Q(x^*)\geq1/2$; their values are the first row below.
Here and below $ij$ means boundaries $\{i,j\}$.
\[
\begin{array}{c|cccccccccc}
Q&12&13&14&15&23&24&25&34&35&45\\\hline
J_Q(x^*)&1&3/4&5/6&1/2&1/2&5/6&1/2&3/4&1/2&1/2\\
\mathbb E_\mu J_Q&2/3&2/3&2/3&2/3&2/3&13/18&5/6&5/6&2/3&2/3
\end{array}
\]
For a fixed internal random seed, every completion agreeing with the
queried answers follows the same adaptive transcript and stopping rule.
Its terminal interval is therefore the final-set interval used above.
Fewer queries cannot decrease width.  Thus every randomized adaptive
policy has expected width at least $1/2$ at $x^*$.

For the batch upper bound, mix $12,23,35$ with probabilities
$1/6,1/3,1/2$.  Its expected joint widths at $e_1,\ldots,e_6$ are
$(2/3,2/3,2/3,2/3,2/3,1/2)$.  Each $J_Q$ is a maximum of linear segment
widths and hence convex; these vertex values bound the expectation
throughout $F$ by $2/3$.  For the reverse bound, place adverse masses
$\mu=(1/3,0,1/6,1/6,1/3,0)$ on those vertices.  The table's second row
shows that every batch has expectation at least $2/3$, and therefore so
does every batch mixture.  Its worst-record expectation is no smaller.
\end{proof}
This is a constructed separation for \emph{mean joint width}, not a
precision-success probability.  It does not follow from the scalar
classification: a maximum of concave local widths need not be concave.
The argument uses established vector recovery and finite-game bounds
\citep{krieg2025}; its claim is the explicit shared-prefix separation.
The two contrasts can be realized with a common forecast baseline
$a=0$ and $b^{(r)}=3c^{(r)}$ in a six-target paired MSE.  Adding other
contrasts would define a different joint objective.
